% Part: lambda-calculus
% Chapter: church-rosser
% Section: parallel-beta-eta-reduction

\documentclass[../../../include/open-logic-section]{subfiles}

\begin{document}

\olfileid[id]{lam}{cr}{pbe}

\olsection{Reduksi-$\beta\eta$ Paralel}

Dalam bagian ini kita membuktikan sifat Church--Rosser bagi
reduksi-$\beta\eta$ paralel, yaitu gagasan reduksi paralel yang bersesuaian
dengan reduksi-$\beta\eta$.

\begin{defn}[Reduksi-$\beta\eta$ paralel, $\beredpar$] \ollabel{defn:beredpar}
  \emph{Reduksi-$\beta\eta$ paralel} ($\beredpar$) pada suku didefinisikan
  secara induktif sebagai berikut:
  \begin{enumerate}
    \item \ollabel{defn:beredpar1} $x \beredpar x$.
    \item \ollabel{defn:beredpar2} Jika $N \beredpar N'$, maka
      $\lambd[x][N]\beredpar\lambd[x][N']$.
    \item \ollabel{defn:beredpar3} Jika $P\beredpar P'$ dan $Q\beredpar Q'$,
      maka $PQ\beredpar P'Q'$.
    \item \ollabel{defn:beredpar4} Jika $N\beredpar N'$ dan $Q\beredpar Q'$,
      maka $(\lambd[x][N])Q\beredpar\Subst{N'}{Q'}{x}$.
    \item \ollabel{defn:beredpar5} Jika $N\beredpar N'$, maka
      $\lambd[x][Nx]\beredpar N'$, dengan syarat $x\notin\FV{N}$.
  \end{enumerate}
\end{defn}

\begin{thm}\ollabel{thm:refl}
  $M \beredpar M$.
\end{thm}

\begin{proof}
  Latihan.
\end{proof}

\begin{prob}
  Buktikan \olref[lam][cr][pbe]{thm:refl}.
\end{prob}

\begin{defn}[Pengembangan lengkap-$\beta\eta$]\ollabel{defn:becd}
  \emph{Pengembangan lengkap-$\beta\eta$}~$\becd{M}$ dari~$M$ didefinisikan
  sebagai berikut:
  \begin{align}
    \becd{x} &= x \ollabel{defn:becd1} \\
    \becd{(\lambd[x][N])} &= \lambd[x][\becd{N}] \ollabel{defn:becd2}\\
    \becd{(PQ)} &= \becd{P}\becd{Q} && \text{jika $P$ bukan abstrak-$\lambd$}
    \ollabel{defn:becd3} \\
    \becd{((\lambd[x][N])Q)} &= \Subst{\becd{N}}{\becd{Q}}{x}
                             \ollabel{defn:becd4} \\
    \becd{(\lambd[x][Nx])} &= \becd{N} \ollabel{defn:becd5} &
    \text{jika $x\notin\FV{N}$}
  \end{align}
\end{defn}

\begin{lem}\ollabel{lem:comp}
  Jika $M\beredpar M'$ dan $R\beredpar R'$, maka $\Subst{M}{R}{y}
  \beredpar\Subst{M'}{R'}{y}$.
\end{lem}

\begin{proof}
  Dengan induksi pada !!{derivation} $M\beredpar M'$.

  Empat kasus pertama tepat sama dengan kasus-kasus dalam
  \olref[pb]{lem:comp}. Jika aturan terakhir adalah
  \olref{defn:beredpar5}, maka $M=\lambd[x][Nx]$ dan $M'=N'$ untuk suatu
  $x$, $N$, dan~$N'$ dengan $x\notin\FV{N}$ serta $N\beredpar N'$. Kita ingin
  menunjukkan
  $\Subst{(\lambd[x][Nx])}{R}{y}\beredpar\Subst{N'}{R'}{y}$, yaitu
  $\lambd[x][\Subst{N}{R}{y}x]\beredpar\Subst{N'}{R'}{y}$. Hal ini mengikuti
  dari \olref{defn:beredpar}\olref{defn:beredpar5} dan hipotesis induksi.
\end{proof}

\begin{lem}\ollabel{lem:cont}
  Jika $M\beredpar M'$, maka $M'\beredpar\becd{M}$.
\end{lem}

\begin{proof}
  Dengan induksi pada !!{derivation} $M\beredpar M'$.

  Empat kasus pertama sama dengan kasus-kasus dalam \olref[pb]{lem:cont}. Jika
  aturan terakhir adalah \olref{defn:beredpar5}, maka
  $M=\lambd[x][Nx]$ dan $M'=N'$ untuk suatu $x$, $N$, dan~$N'$ dengan
  $x\notin\FV{N}$ serta $N\beredpar N'$. Kita ingin menunjukkan
  $N'\beredpar\becd{(\lambd[x][Nx])}$, yaitu $N'\beredpar\becd{N}$, yang
  langsung mengikuti hipotesis induksi.
\end{proof}

\begin{thm}\ollabel{thm:cr}
  $\beredpar$ mempunyai sifat Church--Rosser.
\end{thm}

\begin{proof}
  Langsung dari \olref{lem:cont}.
\end{proof}

\end{document}
